\chapter{2022年「普通物理2」\titlelogo{https://sci.hdu.edu.cn}{HDU}期中试题}
\date{2022年10月25日}{陈许敏}{A0715012s}
{未央物理讲师 \href{https://qm.qq.com/q/UPbGudx8cK}{\textbf{Axia}}}
{\qrcode{https://h5.qun.qq.com/s/AQoh8mQn7O}}{\href{https://h5.qun.qq.com/s/AQoh8mQn7O}{HDU物理营 \!$^{\dagger}$}}
{\qrcode{http://weixin.qq.com/r/mhKYgFTEMQlOrRDj90eI}}{\href{http://weixin.qq.com/r/mhKYgFTEMQlOrRDj90eI}{未央学社公众号}}

\section*{第1部分（每格3分，共30分）}

\begin{problem}[简谐振动]
    一个小球与轻弹簧组成的系统，沿$x$轴做振幅为$A$的简谐振动.若$t=0$时，小球位于$x=\frac{\sqrt{2}}{2}A$处，向$x$负方向运动，则小球简谐运动的初相位是\ans{$\frac{\pi}{4}$}.
\end{problem}
\begin{solution}*
    画出相图后便可轻松判断.
\end{solution}

\begin{problem}[简谐振动的合成]
    一个质点同时参与两个在同一直线上的简谐振动：$x_1=4\cos{\ab(2t+\frac{\pi}{6})},\ x_2=2\cos{\ab(2t-\frac{\pi}{6})}$.该质点合振动的振幅大小为\ans{$2\sqrt{7}$}.
\end{problem}
\begin{solution}*
    $A=\sqrt{A_1^2+A_2^2+2A_1A_2\cos{\Delta\varphi}}=2\sqrt{7}$.
\end{solution}

\begin{problem}[平面简谐波的物理量]
    一束光在某介质材料里的电场方程表达式为$E=\cos{\ab(0.01\pi x-2\pi\times 10^6t)}$ (SI)，则该光波的传播速度是\ans{$2\times 10^{8}\mathrm{m/s}$}，该介质材料的折射率是\ans{$1.5$}.
\end{problem}
\begin{solution}*
电场方程可改写为$E=\cos{\ab[\omega\ab(t-\frac{x}{u})]}$，和题中表达式比对得$0.01\pi x=\frac{2\pi\times 10^6x}{u}$，由此得光波的传播速度$u=2\times 10^{6}\mathrm{m/s}$，介质的折射率$n=\frac{c}{u}=1.5$.
\end{solution}

\begin{problem}[光的偏振]
    自然光和线偏振光混合光束，通过一偏振片时，随着偏振片以光的传播方向为轴的转动，透射光的强度也跟着变化.如最强和最弱的光强之比为$5:1$，那么入射光中自然光和线偏振光的强度之比为\ans{$1:2$}.
\end{problem}
\begin{solution}
    设自然光为$x$，线偏振光为$y$，则最强和最弱的光强之比为
    \[\frac{1/2x+y}{1/2x}=5:1\]

    由此得自然光和线偏振光的强度之比$x:y=1:2$.
\end{solution}

\begin{problem}[波的干涉]
    两列波相遇发生干涉现象的条件是：振动频率相同，振动方向相同和\ans{相位差恒定}.
\end{problem}

\begin{problem}[双缝干涉]
    杨氏双缝干涉实验中，如用折射率$n_1=1.58$得云母片遮盖缝隙$S_1$，用相同厚度折射率$n_2=1.70$的薄玻璃片遮盖缝$S_2$，将使屏幕上原来未放玻璃片时的中央亮纹变为第七级亮纹.已知单色光波长$\lambda=550.0\mathrm{nm}$，则玻璃片的厚度为\ans{$32.08\mu\mathrm{m}$}.
\end{problem}
\begin{solution}*
    光程差的变化为$7\lambda=e\ab(1.70-1)+e\ab(1.58-1)$，由此可得玻璃片的厚度为$32.08\mu\mathrm{m}$.
\end{solution}

\begin{problem}[起偏角]
    玻璃的折射率为$1.50$，当自然光由空气中射向玻璃而反射时，反射光为线偏振光，则入射角\ans{$\arctan{\frac{3}{2}}$}.
\end{problem}
\begin{solution}*
    起偏角$\theta_B=\arctan{\frac{n_2}{n_1}}=\arctan{\frac{3}{2}}$.
\end{solution}

\begin{problem}[多普勒效应]
    一辆高速运动的警车发出警铃，声音频率为$f_s=1400\mathrm{Hz}$，从左向右移动，速度大小为$70\mathrm{m/s}$.一个观察者从右向左移动，速度大小为$10\mathrm{m/s}$.假设空气中没有风，空气中的声速$u=340\mathrm{m/s}$.观察者与警车相遇之前，观察者听到警铃声音的频率是\ans{$1814.8\mathrm{Hz}$}；观察者与警车相遇之后，观察者听到警铃声音的频率是\ans{$1126.8\mathrm{Hz}$}.
\end{problem}
\begin{solution}
    已知多普勒效应观察者 (Observer) 和发射源 (Source) 的频率关系为
    \[\nu=\frac{u\pm v_o}{u\mp v_s}\nu_0\]

    $v_o$为观察者速度，接近为$+$，远离为$-$；$v_s$为发射源速度，接近为$-$，远离为$+$.观察者所听频率分别为
    \[f_1=\frac{340+10}{340-70}\times 1400\mathrm{Hz}\approx 1814.8\mathrm{Hz},\ f_2=\frac{340-10}{340+70}\times 1400\mathrm{Hz}\approx 1126.8\mathrm{Hz}\]
\end{solution}

\section*{第2部分（每题10分，共40分）}
\begin{problem}[简谐振动][10]
    一个质量为$1\mathrm{kg}$的物块沿$x$轴做简谐振动，振幅为$10\mathrm{cm}$，最大速度$2\times 10^{-2}\mathrm{m/s}$.在时间$t=4\mathrm{s}$时，物块位于$5\mathrm{cm}$.求
    \begin{enumerate}
        \item 简谐运动的周期和最大加速度.
        \item 简谐运动的位移方程.
    \end{enumerate}
\end{problem}
\begin{solution}
    \begin{enumerate}
        \item 由$v_{\max}=\omega,\ T=\frac{2\pi}{\omega}$得$T=10\pi\mathrm{s},\ \omega=0.2\mathrm{s}^{-1},\ a_{\max}=\omega^2A=4\times 10^{-3}\mathrm{m/s}^2$.\point4
        \item $t=4\mathrm{s}$时$x=\frac{1}{2}A,\ v<0$，所以此时相位$\varphi=\frac{\pi}{3}$.故位移方程为\point2
        \[x=0.1\cos{\ab[0.2\ab(t-4)]+\frac{\pi}{3}}=0.1\cos{\ab(0.2t+0.079\pi)}\eqno\point4\]
    \end{enumerate}
\end{solution}

\begin{problem}[等厚干涉][10]
    空气中，白光以$i_1=30^{\circ}$入射到玻璃表面，玻璃的厚度$h=0.45\mu\mathrm{m}$，折射率$n=1.50$.在可见光波段内（$400\mathrm{nm}\sim 760\mathrm{nm}$），哪些波长的反射光是干涉相长，哪些波长的反射光是干涉相消？
\end{problem}
\begin{solution}
两束光的光程差$\delta=2nh\cos{i_1}+\frac{\lambda}{2}$，干涉相长、相消条件分别为\point2
\[\begin{cases}
    \delta=k\lambda,&k=2,\ \lambda=467.65\mathrm{nm}\\
    \delta=\frac{2k+1}{2}\lambda,&k=2,\ \lambda=584.57\mathrm{nm}
\end{cases}\eqno\point8\]
\end{solution}

\begin{problem}[牛顿环][10]
    空气中，使用波长为$480\mathrm{nm}$平行单色光观察牛顿环.在反射光中测得某一明环的直径为$4.74\mathrm{mm}$，在它外面第$10$个明环的直径为$7.24\mathrm{mm}$.求
\begin{enumerate}
    \item 平凸透镜的曲率半径.
    \item 直径为$4.74\mathrm{mm}$明环的条纹级数$k$.
    \item 假设把整个装置放入水中($n=1.33$)，原直径为$4.74\mathrm{mm}$明环的新直径.
\end{enumerate}
\end{problem}
\begin{solution}
\begin{enumerate}
    \item 根据明环半径公式
    \[r_k=\sqrt{\frac{2k-1}{2}\lambda R}=2.37\mathrm{mm},\ r_{k+10}=\sqrt{\frac{2\ab(k+10)-1}{2}\lambda R}=3.62\mathrm{mm}\eqno\point2\]
    
    由此得$R=\ab(r_{k+10}^2-r_k^2)/10\lambda=1.56\mathrm{m}$.\point2
    \item 由$r_k=\sqrt{\frac{2k-1}{2}\lambda R}$得$k=\frac{r_k^2}{\lambda R}+\frac{1}{2}=8$.\point3
    \item 由于光在介质中的波长与折射率成反比，所以此时$d_k^{\prime}=\frac{d_k}{\sqrt{n}}=4.11\mathrm{mm}$.\point3
\end{enumerate}
\end{solution}

\begin{problem}[光栅][10]
    波长为$\lambda=600\mathrm{nm}$平行光入射到一个透射光栅，光栅的狭缝宽度$a=1\mu\mathrm{m}$，不透明宽度$b=3\mu\mathrm{m}$.问
\begin{enumerate}
    \item 平行光垂直入射时，在光屏上观察到最高谱线极次是多少.
    \item 平行光垂直入射时，在单缝衍射主极大范围内，可以观测哪些级次的谱线.
    \item 以$i=30^{\circ}$角度入射光栅，在光屏上观察到最高谱线级次是多少.
    \item 以$i=30^{\circ}$角度入射光栅，在光屏上可以看到哪些级次的谱线.
\end{enumerate}
\end{problem}
\begin{solution}
\begin{enumerate}
    \item 由光栅方程
    \[\ab(a+b)\sin{\theta}=k\lambda\eqno\point1\]
    
    当$\sin{\theta}=1$时，可得最大级次
    \[k_{\max}=\left\lfloor\frac{a+b}{\lambda}\right\rfloor=6\eqno\point1\]
    \item 利用缺级条件$a\sin{\theta}=k^{\prime}\lambda$得$\frac{k^{\prime}}{k}=\pm\frac{a+b}{a}=\pm 4,\pm 8,\cdots$.所以可见级次谱线为$0,\pm 1,\pm 2,\pm 3,\pm 5,\pm 6$.\point1
    \item 由光栅方程
    \[\ab(a+b)\ab(\sin{\theta}-\sin{i})=k\lambda\eqno\point2\]

    令$\theta=\pm\frac{\pi}{2}$，得两侧最高级次分别为
    \[k_{+\pi/2}=\left\lfloor\frac{a+b}{2\lambda}\right\rfloor=3,\ k_{-\pi/2}=-\frac{3\ab(a+b)}{2\lambda}=-10\eqno\point2\]

    由于$-\frac{\pi}{2}$处在现实中无法观测，所以两侧的最高级次分别为$-3$和$9$.\point1
    \item 由于$\pm 4,\pm 8,\cdots$缺级，所以可见级次谱线为$0,\pm 1,\pm 2,\pm 3,-5,-6,-7,-9$.\point2
\end{enumerate}
\end{solution}

\newpage
\section*{第3部分（每题15分，共30分）}
\begin{problem}[驻波][15]
    入射波的波函数$y=0.01\cos{\ab(2\pi t-\pi x+\frac{\pi}{6})}$，在$x_0=\frac{9}{2}$处发生反射，反射点为节点.求
\begin{enumerate}
    \item 反射波的波函数.
    \item 合成驻波的波函数.
    \item 坐标原点和$x_0=\frac{9}{2}$之间，各波腹和波节的位置.
\end{enumerate}
\end{problem}
\begin{solution}
\begin{enumerate}
    \item 反射后会有$\pi$相位的半波损失，所以此时$x=\frac{9}{2}$处的反射波振动方程为
    \[y_2\ab(\frac{9}{2},t)=0.01\cos{\ab(2\pi t-\pi\times\frac{9}{2}+\frac{\pi}{6}+\pi)}=0.01\cos{\ab(2\pi t-\frac{10\pi}{3})}\eqno\point2\]

    由此可得反射波的波函数
    \[y_2\ab(x,t)=0.01\cos{\ab(2\pi t+\pi\ab(x-\frac{9}{2})-\frac{10\pi}{3})}=0.01\cos{\ab(2\pi t+\pi x+\frac{\pi}{6})}\eqno\point3\]
    \item 驻波的波函数
    \[y=y_1+y_2=0.02\cos{\ab(2\pi t+\frac{\pi}{6})}\cos{\ab(\pi x)}\eqno\point6\]
    \item 令$y=\pm A$，得波腹的位置$x=0,1,2,3,4$.\point2
    \item 令$y=0$，得波节的位置$x=\frac{1}{2},\frac{3}{2},\frac{5}{2},\frac{7}{2},\frac{9}{2}$.\point2
\end{enumerate}
\end{solution}

\begin{problem}[弹簧振子][15]
    一个质量为$m=1\mathrm{kg}$的盘子刚性连接竖直悬挂的轻弹簧下端，弹簧的劲度系数为$k=90\mathrm{N/m}$.盘子在竖直方向做简谐运动，振幅$A=10.0\mathrm{cm}$.现有一个质量$m_2=1\mathrm{kg}$的物体自由落下掉在盘上，没有反弹.当\textbf{盘子位于向上最大位移处}，盘子与物体发生碰撞，物体从开始自由自由落体发生处，距离为$h=20\mathrm{cm}$，假设碰撞瞬间完成.求碰撞后盘子和物体组成的系统，它的振动周期、振幅和振动能量.
\end{problem}
\begin{solution}
\begin{itemize}
    \item 系统的振动周期$T=2\pi\sqrt{\frac{m_1+m_2}{k}}=0.94\mathrm{s}$，角频率$\omega=\frac{2\pi}{T}=6.71\mathrm{s}^{-1}$.\point3
    \item 物体自由落体到盘上速度$v_0=\sqrt{2gh}=2\mathrm{m/s}$，利用与盘子完全非弹性碰撞中动量守恒得二者初速度\point1
    \[m_2v_0=\ab(m_1+m_2)v_1,\ v_1=\frac{m_2}{m_1+m_2}v_0=1\mathrm{m/s}\eqno\point2\]
    \item 碰撞后，盘子和物块新的平衡点较盘子原来的平衡点下移$\Delta x=\frac{m_2g}{k}=\frac{1}{9}\mathrm{m}$，此时盘子与物块距离二者新平衡点距离$x=A+\Delta x=0.21\mathrm{m}$.由旋转矢量法，此时系统的振幅\point3
    \[A^{\prime}=\sqrt{x^2+\ab(\frac{v_1}{\omega})^2}=25.84\mathrm{cm}\eqno\point3\]

    振动能量
    \[E=\frac{1}{2}kA^{\prime^2}=3.01\mathrm{J}\eqno\point3\]
\end{itemize}
\end{solution}